\documentclass{beamer}
\usetheme{Madrid}
\useoutertheme[subsection=false]{smoothbars}
\usecolortheme{default}

\definecolor{smoothbarblue}{RGB}{38,38,134}
\setbeamercolor{section in head/foot}{fg=white,bg=smoothbarblue}
\usepackage{amssymb}

\theoremstyle{definition}
\newtheorem{mytheor}[theorem]{Theorem}
\newtheorem{mydef}[theorem]{Definition}
\newtheorem{proposition}[theorem]{Proposition}
\newtheorem{remark}[theorem]{Remark}
\newtheorem{myexample}[theorem]{Example}

\setbeamertemplate{navigation symbols}{}

\title[Stochastic optimal control]{Stochastic optimal control}
\author[Narek Sargsian, Ivan Khromyshchenko]{Narek Sargsian, Ivan Khromyshchenko}
\institute[HSE FES]{HSE FES Probability Theory Club}
\date{28 February 2026}
\usepackage{graphicx}

\usepackage{amsmath}

\newcommand{\R}{\mathbb{R}}
\newcommand{\E}{\mathbb{E}}
\renewcommand{\P}{\mathbb{P}}
\DeclareMathOperator{\Var}{Var}
\DeclareMathOperator{\Cov}{Cov}
\DeclareMathOperator{\Corr}{Corr}
\DeclareMathOperator{\sign}{sign}
\DeclareMathOperator{\rank}{rank}
\DeclareMathOperator{\plim}{plim}

\usepackage{euscript}

\setbeamertemplate{footline}{%
  \leavevmode%
  \hbox{%
    \begin{beamercolorbox}[wd=.25\paperwidth,ht=2.25ex,dp=1ex,center]{title in head/foot}%
      \insertshorttitle
    \end{beamercolorbox}%
    \begin{beamercolorbox}[wd=.50\paperwidth,ht=2.25ex,dp=1ex,center]{author in head/foot}%
\insertshortauthor\hspace*{0.6em}(\insertshortinstitute)
    \end{beamercolorbox}%
    \begin{beamercolorbox}[wd=.25\paperwidth,ht=2.25ex,dp=1ex]{date in head/foot}%
      \hbox to .25\paperwidth{%
        \hbox to .18\paperwidth{\hfil\insertshortdate\hfil}%
        \hbox to .05\paperwidth{\hfil\insertframenumber/\inserttotalframenumber\hspace*{1ex}}%
      }%
    \end{beamercolorbox}%
  }%
}

\begin{document}

\begin{frame}
    \titlepage
\end{frame}

\section{Introduction}
\subsection{}

\begin{frame}{Introduction}

\begin{itemize}
    \item Motivation for the use of stochastic optimal control
    \item Tasks that can be solved using stochastic optimal control
    \item Necessary tools
    \item An optimal periodic dividend and risk control problem for an insurance company
    \item Merton's Optimal investment problem  
    \item Next steps
\end{itemize}
    
\end{frame}

\begin{frame}{Motivation}

Deterministic optimal control assumes the system evolves predictably. In practice, however, many processes are stochastic. Incorporating randomness into the control framework therefore makes models more realistic.

\end{frame}

\begin{frame}{Problem scope overview}

\begin{itemize}
    \item Insurance company problem
    \item Merton's Optimal investment problem
    \item Passport options
    \item Avellaneda-Stoikov model
\end{itemize}
\end{frame}


\section{Math}
\subsection{}

\begin{frame}{HJB}

Define the value function (cost minimization):
\begin{equation}
    V(t, x) = \min_{u(\cdot)}\left[\int_t^\tau C(s, x(s), u(s)) \, ds + D(x(\tau)) \right]
\end{equation}

subject to the \textbf{state dynamics}:
\begin{equation}
    \dot{x}(s) = f(s, x(s), u(s)), \qquad x(t) = x.
\end{equation}

where:
\begin{itemize}
    \item $C(t,x,u)$ is the \textbf{running cost} (cost rate over time),
    \item $D(x)$ is the \textbf{terminal cost} at the horizon $\tau$.
\end{itemize}

\end{frame}

\begin{frame}{HJB}
    The HJB equation is
    \begin{equation}
    \partial_t V(t,x)
    + \min_{u} \left\{
        \partial_x V(t,x)\, f(t,x,u)
        + C(t,x,u)
    \right\} = 0,
    \end{equation}
    with terminal condition
    \begin{equation}
        V(\tau,x) = D(x).
    \end{equation}
\end{frame}





\begin{frame}{Stochastic HJB}

In the stochastic case, define the value function (cost minimization):
\begin{equation}
    V(t,x) = \min_{u(\cdot)} \E_{t,x}\!\left[\int_t^\tau C(s, X_s, u_s) \, ds + D(X_\tau) \right],
\end{equation}
where \(\E_{t,x}\) denotes conditional expectation given \(X_t=x\).

The controlled state process satisfies
\begin{equation}
    dX_s = f(s, X_s, u_s)\,ds + \sigma(s, X_s, u_s)\,dW_s, \qquad X_t=x.
\end{equation}

\end{frame}

\begin{frame}{Stochastic HJB}
    The HJB equation is
    \begin{equation}
        \partial_t V(t,x) + \min_{u} \left\{ \mathcal{L}^u V(t,x) + C(t,x,u) \right\} = 0,
    \end{equation}
    with terminal condition \(V(\tau,x)=D(x)\).

    where

    \begin{equation*}
        \mathcal{L}^u V(t,x) = \sum_{i=1}^{d} b_i(t,x,u)\,\partial_{x_i} V(t,x)
    + \frac{1}{2} \sum_{i,j=1}^{d} a_{ij}(t,x,u)\,\partial^{2}_{x_i x_j} V(t,x),
    \end{equation*}
    where \(b(t,x,u)=f(t,x,u)\) and \(a(t,x,u)=\sigma(t,x,u)\sigma(t,x,u)^\top\).
    
\end{frame}

\begin{frame}{Transition to stationarity}
    If the problem is time-invariant (time-homogeneous dynamics and costs), the value function does not depend explicitly on time:
\[
V(t,x) = V(x).
\]
Therefore,
\[
\partial_t V(t,x) = 0,
\]
and the HJB equation becomes stationary.
\end{frame}


\begin{frame}{Algorithm}
    \begin{itemize}
        \item Defining the dynamics of the process (SDE)
        \item Defining the objective functional and value function
        \item Admissible control set
        \item Hamilton--Jacobi--Bellman (HJB) equation
        \item First-order conditions and optimal control
    \end{itemize}
\end{frame}

\section{Merton's Portfolio Problem}
\subsection{}

\begin{frame}{Motivation}
    Before Merton, portfolio theory relied on the Markowitz model. The latter calculates weights only once. Merton's model chooses a dynamic policy.

    \textbf{Markowitz's Model in the utility form:}
    \[
    \max_{w \in \mathbb{R}^n}\; w^\top \mu - \frac{\gamma}{2}\, w^\top \Sigma w
    \quad \text{s.t.} \quad
    \mathbf{1}^\top w = 1
    \]

\begin{itemize}
    \item $n$ is the number of risky assets.
    \item $R \in \mathbb{R}^n$ is the random vector of single-period risky-asset returns.
    \item $w \in \mathbb{R}^n$ is the vector of portfolio weights in the risky assets.
    \item $\mu \in \mathbb{R}^n$ is the vector of expected returns, $\mu = \E[R]$.
    \item $\Sigma \in \mathbb{R}^{n \times n}$ is the covariance matrix of returns, $\Sigma = \Cov(R)$.
    \item $w^\top \mu$ is the portfolio expected return.
    \item $w^\top \Sigma w$ is the portfolio return variance.
    \item $\gamma > 0$ is the risk-aversion parameter (larger $\gamma$ puts more penalty on variance).
\end{itemize}
\end{frame}

\begin{frame}{Informal Problem Statement}
The investor chooses how much to consume over time $c_t$ and how to allocate wealth between risky assets and the risk-free asset.
\end{frame}

\begin{frame}{Problem Notation}
    The utility function is of the constant relative risk aversion (CRRA) form (Merton's parameterization):
    \begin{equation}
    u(x)=\frac{x^{\gamma}}{\gamma}, \qquad \gamma<1,\ \gamma\neq 0.
    \end{equation}
    (Then relative risk aversion is $1-\gamma$.)
    
    A utility function other than CRRA can be used.
    
    \begin{itemize}
        \item \(\mathcal{A}\) is the set of strategies $a = (\pi, c)$ and \(\tau_a\) is the stopping time.\\ 
        \item $X_t$ is the wealth at time $t$
        \item The initial wealth is $x$
    \end{itemize}
\end{frame}

\begin{frame}{Problem Statement}

The objective is
\begin{equation}
J^{T}(x;\pi,c)=\E\!\left[\int_{0}^{T} e^{-\rho s}\,u(c_s)\,ds \;+\; e^{-\rho T}\,u(X_{T})\bigg| \ X_0 = x\right]
\end{equation}\\

where the wealth dynamics follow the SDE

\begin{equation}
dX_t = \left[(r + \pi_t(\mu - r))X_t - c_t\right]dt + X_t \pi_t \sigma dW_t
\end{equation}\\

\begin{itemize}

    \item \(\rho > 0\) is the subjective discount rate
    \item \(\mu\) is the expected return of the risky asset
    \item \(\sigma\) is the volatility of the risky asset
    \item \(r\) is the risk-free rate
\end{itemize}
\end{frame}

\begin{frame}{Defining the value function}

Assuming \(T \to \infty\) (and a transversality condition so the terminal term vanishes), define the infinite-horizon objective
\begin{equation}
J(x;\pi,c)= \E \left[\int_{0}^{\infty} e^{-\rho s}\,u(c_s)\,ds \bigg| \ X_0 = x\right].
\end{equation}

The value function is
\begin{equation}
V(x) = \sup_{(\pi, c) \in \mathcal{A}} J(x;\pi,c).
\end{equation}

\end{frame}


\begin{frame}{Stationarity (infinite horizon)}
With time-homogeneous coefficients and an infinite horizon, the value function is stationary: \(V(t,x)=V(x)\).
The HJB equation therefore reduces to an ODE in wealth \(x\).
\end{frame}

\begin{frame}{HJB equation}
    The optimality equation becomes an ODE in wealth:
    \begin{equation}
    0=\max_{c,\pi}
    \Big[u(c)-\rho V(x)
    +V'(x)\big((\pi(\mu-r)+r)x-c\big)
    +\frac12 V''(x)\sigma^2\pi^2 x^2\Big].
    \end{equation}
\end{frame}

\begin{frame}{Solution}
    The first-order conditions are
    \begin{equation}
    0=u'(c)-V'(x),
    \end{equation}
    \begin{equation}
    0=(\mu-r)V'(x)+V''(x)\sigma^2\pi x.
    \end{equation}
\end{frame}

\begin{frame}{Solution}
    The optimal decision rules are:
    
    \begin{equation}
    \boxed{\ \pi_\infty^*=\frac{\mu-r}{\sigma^2(1-\gamma)}\ }
    \end{equation}
    
    \begin{equation}
    \boxed{\ c_t^*=
    \left\{
    \frac{\rho}{1-\gamma}
    -\gamma\left[
    \frac{(\mu-r)^2}{2\sigma^2(1-\gamma)^2}
    +\frac{r}{1-\gamma}
    \right]
    \right\}X_t\ }
    \end{equation}
\end{frame}



\section{Insurance}
\subsection{}

\begin{frame}{Motivation}

The surplus of an insurance company can be modeled as a stochastic process.
This allows us to apply stochastic optimal control methods to the problem of managing the reinsurance part and dividend payments.

\end{frame}

\begin{frame}{Problem Statement}
The surplus process can be described as follows

\begin{equation}
    X_t = x_0 + \lambda a [1 + \eta] t - \lambda a [1 - u_t] [1 + \mu] t - u_t \sum_{j=1}^{N_t^{\lambda}} Y_j - L_t, \quad \text{for } t \ge 0,
    \label{eq:surplus_process}
\end{equation}

\begin{itemize}
    \item $x_0$ denotes the initial surplus of the insurance company.
    
    \item $u_t$ represents the proportion of risk transferred to the reinsurer at time $t$.
    
    \item $Y_j$ is the size of the $j$-th insurance claim, where $\E[Y] = a$ and $\Var(Y) = \widehat{\sigma}^2$.
    
    \item $\lambda$ is the intensity (arrival rate) of insurance claims.
    
    \item $\eta$ denotes the premium inflow rate from policyholders.
    
    \item $\mu$ is the payment rate from the insurer to the reinsurance company.
    
    \item $L_t$ is a non-decreasing process representing the cumulative dividends paid up to time $t$.
\end{itemize}



\end{frame}

\begin{frame}{Problem Statement}
We can rewrite the process as

\[
X_t = x_0 + \lambda a \bigl[ \eta - [1 - u_t] \mu \bigr] t + u_t 
    \underbrace{\left[ \lambda a t - \sum_{j=1}^{N_t^\lambda} Y_j \right]}_{Z_t} - L_t, 
    \quad \text{for } t \ge 0.
\]

We have \(\E[Z_t]=0\) and \(\Var(Z_t)=\lambda t(\widehat{\sigma}^2+a^2)\approx \widehat{\sigma}^2\lambda t\). For \(\lambda \gg 1\), a diffusion approximation yields \(Z_t \approx \sigma W_t\), where \(W_t\) is a standard Brownian motion and \(\sigma = \widehat{\sigma}\sqrt{\lambda}\).
\end{frame}

\begin{frame}{Dynamics of the process and defining the dividend payment}
The diffusion approximation is

\begin{equation}
    dX_t = [\eta - [1 - u_t]\mu] dt + \sigma u_t dW_t - dL_t, \quad \text{for } t \ge 0 \text{ and } X_0 = x_0.
    \label{eq:diff_approx}
\end{equation}

Dividend payments may be described as follows
    \begin{equation}
    L_t = \int_0^t \nu_s \, \mathrm{d}N_s^\gamma, \quad \text{for } t \ge 0,
    \label{eq:divid}
    \end{equation}

where \(\gamma > 0\) is the intensity.

\end{frame}

\begin{frame}{Defining the set of controls}
Let $\mathcal{A}$ denote the set of strategies $\pi = (u^{\pi}, L^{\pi})$ such that $u^{\pi} \in [0,1]$, $L^{\pi}$ satisfies \eqref{eq:divid}, and $0 \le \nu_t^{\pi} \le X_{t-}^{\pi}$. 
A strategy $\pi$ is said to be admissible if $\pi \in \mathcal{A}$.

\end{frame}

\begin{frame}{Defining NPV}
For each \(\pi \in \mathcal{A}\), the expected NPV is given by

\begin{equation}
    J^\pi(x) = \E_x\left[\int_0^{\tau^\pi} e^{-\delta t} dL^\pi_t \bigg| X_0 = x \right] = \E_x\left[\int_0^{\tau^\pi} e^{-\delta t} \nu_t^\pi dN^\gamma_t\bigg| X_0 = x\right]
    \label{eq:npv}
\end{equation}

Here, $\E_x$ denotes the conditional expectation given the initial surplus level $x$, 
$\delta > 0$ is the discount rate, and 
$\tau^{\pi} := \inf\{ t > 0 : X_t^{\pi} < 0 \}$ represents the time of ruin. 
Our objective is to determine the value function $V$ defined by

\begin{equation}
V(x) := \sup_{\pi \in \mathcal{A}} J^{\pi}(x), \qquad x \ge 0.
\label{eq:NPV_max}
\end{equation}

There is an assumption that \(V \in C^2(0, \infty)\) and concave.

\end{frame}

\begin{frame}{HJB equation}
To find the optimal \(V\), we construct the HJB equation as follows:

\begin{equation}
    V(x) = \underbrace{(\nu + \E_{x}[V(X_{t} - \nu)]) \gamma t e^{-\delta t}}_{\text{in case of dividend payment}} + \underbrace{(1 - \gamma t) e^{-\delta t} \E_{x}[V(X_{t})]}_{\text{in case no dividend is paid}} + o(t).
    \label{eq:HJB_const}
\end{equation}

\[
(\nu + \E_{x}[V(X_{t} - \nu) - V(x)]) \gamma t e^{-\delta t} + (1 - \gamma t) e^{-\delta t} \E_{x}[V(X_{t}) - V(x)] = o(t)
\]

\noindent
\textbf{Intuition:} The current value equals the expected discounted value over a small time interval $t$ plus a possible dividend payment. In addition, \(\gamma t\) comes from the first-order expansion of the Poisson probabilities (up to \(o(t)\)).

\end{frame}

\begin{frame}{HJB equation}
We find the change over time of the discounted part and the value function by differentiating \(Y_t = e^{-\delta t}V(X_t)\).

\begin{equation}
    dY_t = e^{-\delta t}dV(X_t) - \delta e^{-\delta t} V(X_t) \Rightarrow{} dY_t = e^{- \delta t} (dV(X_t) - \delta V(X_t))
    \label{eq:diff_of_Y_t}
\end{equation}

For \(dV(X_t)\) we use It\^{o}'s formula. The dynamics are the same as in \eqref{eq:diff_approx}, but without the dividend term \(L_t\). Set \(a_s := \eta - (1-u_s)\mu\).
    
\end{frame}

\begin{frame}{Derivation of HJB equation}
We consider two cases:

\begin{enumerate}
    \item The company does not pay a dividend.
    \item The company pays a dividend.
\end{enumerate}

Consider the first case:

\begin{align}
    &e^{-\delta t}V(X_t) - V(x) = \int_0^t e^{-\delta s}
        [a_s V'(X_s)
        + \frac{1}{2}\sigma^2 u_s^2 V''(X_s) \notag \\
    & \quad - \delta V(X_s)]ds 
     + \int_0^t e^{-\delta s} \sigma u_s V'(X_s)\, dW_s.
    \label{eq:first_case_1}
\end{align}
    
\end{frame}

\begin{frame}{Derivation of HJB equation}

Take expectations in \eqref{eq:first_case_1}.

\begin{align}
    &\E_x\left[e^{-\delta t}V(X_t) - V(x)\right] = \notag \\
    & \quad = \E_x\left[\int_0^t e^{-\delta s}
        \left[a_s V'(X_s)
        + \frac{1}{2}\sigma^2 u_s^2 V''(X_s)
        - \delta V(X_s)\right]ds\right]
    \label{eq:first_case_2}
\end{align}

Let's do the same for the second case.

\end{frame}

\begin{frame}{Derivation of HJB equation}
\begin{align}
e^{-\delta t}V(X_{t-} - \nu) - V(x)
&= \int_0^t e^{-\delta s}
\Big[
a_s V'(X_s)
+ \tfrac{1}{2}\sigma^2 u_s^2 V''(X_s)
- \delta V(X_s)
\Big] ds \notag \\
&\quad + \int_0^t e^{-\delta s}
\sigma u_s V'(X_s)\, dW_s \notag \\
&\quad + e^{-\delta t}
\Big(
V(X_{t-}-\nu)-V(X_{t-})
\Big).
\label{eq:second_case_1}
\end{align}
Taking expectations yields
\begin{align}
&\E_x \!\left[e^{-\delta t}V(X_{t-}-\nu)-V(x)\right] = \notag \\
&= \E_x \!\left[\int_0^t e^{-\delta s}
    \left[a_s V'(X_s) + \frac{1}{2} \sigma^2 u_s^2 V''(X_s) - \delta V(X_s)\right] ds \right] + \notag \\
&\quad + \E_x\left[e^{-\delta t}\left[ V(X_{t-}-\nu) - V(X_{t-})\right] \right].
\label{eq:second_case_2}
\end{align}
\end{frame}

\begin{frame}{Derivation of HJB equation}
Substituting \eqref{eq:first_case_2} and \eqref{eq:second_case_2} into \eqref{eq:HJB_const},

\begin{align}
    &\nu\gamma t e^{-\delta t} + \E_x \!\left[\int_0^t e^{-\delta s}
    \left[a_s V'(X_s) + \frac{1}{2} \sigma^2 u^2_s V''(X_s) - \delta V(X_s)\right] ds \right]\gamma t + \notag \\
    &\quad + \E_x \left[e^{-\delta t}\left[ V(X_{t-}-\nu) - V(X_{t-})\right]\right] \gamma t +  \notag \\
    & + (1 - \gamma t) \E_x\left[\int_0^t e^{-\delta s}
        \left[a_s V'(X_s)
        + \frac{1}{2}\sigma^2 u_s^2 V''(X_s)
        - \delta V(X_s)\right]ds\right] = 0
    \label{eq:done}
\end{align}

\end{frame}

\begin{frame}{Derivation of HJB equation}
Dividing \eqref{eq:done} by \(t\) and letting \(t \to 0\) yields the stationary HJB equation:

\begin{equation}
    \max_{0 \leq \nu \leq x, \ 0 \leq u \leq 1} \underbrace{\mathcal{L}^u V(x) + \gamma \left[V(x - \nu) - V(x) + \nu  \right]}_{\Xi_{x}(u,\nu)} = 0,
    \label{eq:HJB}
\end{equation}
where \(V(0) = 0\) and
\begin{equation*}
    \mathcal{L}^u V(x) = \left[[\eta - [1-u]\mu]V'(x)+ \frac{1}{2}\sigma^2 u^2 V''(x)- \delta V(x)\right].
\end{equation*}

    
\end{frame}

\begin{frame}{Defining the barrier strategy}

For each \( \mathbf{b} = (b_1, b_2)\), define a periodic-classical barrier strategy \(\pi^{\mathbf{b}} = (u^{\mathbf{b}}, L^{\mathbf{b}})\), where \(u^{*}(\cdot)\) is the feedback control on \((0,b_1)\) (specified below):

\begin{equation}
    u_t^{\mathbf{b}} := 
    \begin{cases}
        u^{*}(X_t^{\mathbf{b}}), & \text{if } X_t^{\mathbf{b}} \in (0,b_1), \\
        1, & \text{if } X_t^{\mathbf{b}} \in [b_1,\infty),
    \end{cases}
    \label{eq:barrier_strat}
\end{equation}
\[
L_t^{\mathbf{b}} = \int_{0}^{t} \nu_s^{\mathbf{b}} \, dN_s^\gamma,
\qquad
\nu_s^{\mathbf{b}} = \max\{ X_s^{\mathbf{b}} - b_2, 0 \}.
\]

\end{frame}

\begin{frame}{Defining the barriers}
    The optimal barriers are obtained by solving \eqref{eq:HJB}

    \begin{equation}
    b_1^{\gamma} = \inf \left\{ x > 0 : -\frac{\mu}{\sigma^2} f(x;v) \geq 1 \right\}, \quad
    b_2^{\gamma} = \inf \left\{ x > 0 : v'(x) < 1 \right\},
    \label{eq:optim_bar}
\end{equation}


where \(f(x; v) := \frac{v'(x)}{v''(x)}\) and \(\inf \varnothing = \infty\).

\end{frame}

\begin{frame}{Optimal strategies}
Based on the optimal barriers, the optimal strategy is conjectured to take the form
    \[
    u^{*}(x) =
    \begin{cases}
        -\dfrac{\mu}{\sigma^{2}} f(x;v), & \text{if } x \in (0,b_{1}^{\gamma}),\\[4pt]
        1, & \text{if } x \in [b_{1}^{\gamma},\infty),
    \end{cases}
    \]
    \\
    \[
    \nu^{*}(x) =
    \begin{cases}
        0, & \text{if } x \in (0,b_{2}^{\gamma}),\\[4pt]
        x - b_{2}^{\gamma}, & \text{if } x \in [b_{2}^{\gamma},\infty),
    \end{cases}
    \]
\end{frame}



\begin{frame}{Illustration of the strategies}

\begin{figure}
\centering
\includegraphics[width=\linewidth,height=0.8\textheight,keepaspectratio]{illustratoin.png}
\caption{Illustration of the barriers}
\end{figure}

\end{frame}

\begin{frame}{System of the differential equations for \(v(x)\)}

\begin{align}
-\frac{\mu^2 [v'(x)]^2}{2\sigma^2 v''(x)}
+ [\eta-\mu]v'(x) - \delta v(x) +
& \notag \\
\qquad + \gamma\{x-b_2+v(b_2)-v(x)\}\,\mathbb{1}_{\{x-b_2>0\}}
&= 0 \quad \text{for } x\in(0,b_1),
\end{align}

\begin{align}
\frac{1}{2}\sigma^2 v''(x) + \eta v'(x) - \delta v(x) +
& \notag \\
\qquad + \gamma\{x-b_2+v(b_2)-v(x)\}\,\mathbb{1}_{\{x-b_2>0\}}
&= 0 \quad \text{for } x\in(b_1,\infty),
\end{align}

\[
\text{s.t. } v(0)=0.
\]

\end{frame}


\begin{frame}{Solution for the very expensive case}

\[
v_{\gamma,b}(x)=
\begin{cases}
c_{1,1}(b_2)\,h_1(x), \ \ \ \ x\in(0,b_2),\\[4pt]
c_{1,2}(b_2)e^{\lambda_\gamma(x-b_2)}
\dfrac{\gamma}{\gamma+\delta}\!\left[x-b_2+v_{\gamma,b}(b_2)\right]
\dfrac{\eta}{\gamma+\delta}, \ \ \ \  x\in[b_2,\infty),
\end{cases}
\]\\


where \(\lambda_\gamma\) is the negative root of 
\[
\frac{\sigma^2}{2}r^2+\eta r-(\delta+\gamma)=0,
\]

\[
h_1(x)=e^{\theta_+ x}-e^{\theta_- x}\quad \text{for } x\in(0,\infty),
\]


with \(\theta_+,\theta_-\) the positive and negative root, respectively, of 
\(\frac{\sigma^2}{2}r^2+\eta r-\delta=0.
\)


\end{frame}

\begin{frame}{Solution to the Cauchy problem}
\[
c_{1,1}(b_2)=
\frac{\frac{\gamma}{\gamma+\delta}\!\left[1-\frac{\eta\lambda_\gamma}{\gamma+\delta}\right]}
{h_1'(b_2)-\frac{\delta\lambda_\gamma}{\gamma+\delta}h_1(b_2)}\]
\\
\[
c_{1,2}(b_2)=
\frac{1}{\gamma+\delta}
\left[\delta c_{1,1}(b_2)h_1(b_2)-\frac{\gamma\eta}{\gamma+\delta}\right]
\]

\end{frame}

\section{Conclusion}
\subsection{}

\begin{frame}{Assumptions/Limitations}
\begin{itemize}
    \item In some cases solving the non-linear differential equation analytically is impossible.
    \item For some processes it is quite hard to solve the optimization problem (i.e. compound Poisson process).
    \item Formulation of HJB equation requires rigorous math.
    \item Non-linear differential equations can be solved with PINNs (Physics-informed neural networks)
\end{itemize}

\end{frame}

\begin{frame}{References}
\begin{enumerate}
    \item Karatzas, I., \& Shreve, S. E. (1991). Brownian Motion and Stochastic Calculus (2nd ed.). Springer. (Graduate Texts in Mathematics, Vol. 113).
    \item  Merton, Robert C. (1969). Lifetime Portfolio Selection under Uncertainty: The Continuous-Time Case. The Review of Economics and Statistics, 51(3), 247–257.
    \item Kelbert, M., \& Moreno-Franco, H. A. (2025). An optimal periodic dividend and risk control problem for an insurance company. Insurance: Mathematics and Economics, 125, 103154.
\end{enumerate}
\end{frame}   

\end{document}